\begin{afterword}
\summaryhead

The post opens the sequence on quantum mechanics by saying how it will teach. The author, who is not a physicist but has explained Bayes's theorem, holds that confusion is in our models, not in the world, and that a teacher should not leave students confused. So the sequence will not say that quantum mechanics is supposed to be confusing, or that no one understands it. It will speak of quantum mechanics as normal and make fun of the intuitions that find it strange.

It will also skip the historical story in which matter acts sometimes like billiard balls and sometimes like waves. An electron is neither. The order of discovery need not be the order of teaching, and the sequence will start from the quantum level rather than from experimental results, which the author counts as complicated, high-level phenomena.

Finally, it takes a strictly realist view: the equations describe the territory, and the classical world exists only implicitly within the quantum one. A sizable community of scientists disputes this, the author warns, and non-realist views will be discussed later.

\responsehead

The post is a plan for the sequence, and most of it is announcement. Two things in it deserve credit. The author says plainly that the author is not a physicist, and the post warns readers that a sizable community of scientists disputes the view they are about to be taught.

The main problem is that the post states as settled a question that is still disputed. No one understood quantum mechanics once, it says, but ``we no longer live in that era''; the confusion was ``finally sorted out in the second half of the twentieth century''; the classical world is ``strictly implicit'' in the quantum one; and this is, ``at a core level,'' ``the truth.'' Later posts show that the resolution meant is many-worlds together with decoherence, the work of Everett, Zeh and Zurek, whom this post does not name. The Stanford Encyclopedia reports no consensus among physicists or philosophers on what the theory tells us about the world, and says that decoherence alone does not solve the measurement problem, the question of how a measurement comes to have a single outcome. The author's own ``Many Worlds, One Best Guess,'' a month later, says: ``There is no known reason for the Born probabilities'' (the rule that turns amplitudes into probabilities).

The warning near the end of the post helps, but it comes after these claims, and it describes the dispute as realism against non-realism. Pilot-wave theory, a realist rival in which particles have definite paths, is left out, here and in the rest of the book. The claim that the author's view is ``probably the majority view among theoretical physicists'' is hedged and unsupported. The polls since, of other groups, did not find a majority for a purely realist view.

The teaching argument is reasonable but untested. ``The result of the standard approach is standard confusion'' comes with no evidence about students. The stance the post offers as a departure from the standard introduction, that an electron is neither particle nor wave and must be described on its own terms, is the one Feynman took in the 1964 lecture from which the post takes ``nobody understands quantum mechanics.''

In short: a plan to teach quantum mechanics as normal, open about the author not being a physicist and about dissent, which nonetheless presents one contested interpretation as the settled outcome of history.
\end{afterword}
